\chapter{Learning the Kernel, and What Convergence Costs}\label{ch:em-kernel}

Recovering two numbers is one problem; recovering a whole transition law is
another. This chapter reports what happens when the innovation is represented
nonparametrically, and the finding is not the recovery itself but its
\emph{schedule}: the correlation and the innovation shape converge at very
different rates, and a stopping rule read off the first certifies a kernel that
has not converged in the second. Every non-Gaussian claim in this thesis
depends on the coordinate that settles last, so this is a prerequisite for
Chapter~\ref{ch:em-results} rather than a diagnostic beside it.

\section{What the channel costs the E-step}\label{sec:em-convergence}

What this chapter reports is not a theorem but a measurement, and it is the
one that most directly qualifies the results of Chapter~\ref{ch:em-results}:
how many EM iterations the algorithm of Section~\ref{sec:em-algorithm}
actually needs.

Table~\ref{tab:em-convergence} shows what the channel costs, in two parts: the
through-channel fit at four updates, and where a clean and a channel arm each
finish.

%% GENERATED by tools/make_convergence_numbers.py -- do not hand-edit.
%% Trace: outputs/frozen/exp_27_seed16/shape_trace.csv (16 seeds). Endpoints: outputs/frozen/exp_06/clean_vs_noised_shape.csv.

\begin{center}
\begin{minipage}{\linewidth}\centering
\captionof{table}[What the channel costs in convergence rate]{What the
observation channel costs in rate rather than in where the fit lands. Above:
the through-channel fit at $\nseq = 512$, $C = 8$, as a
median over $16$ seeds at four updates ---
the correlation is essentially in place by $30$ while the innovation shape is
still moving at $120$. Below: where the two arms finish, fit on clean transition
pairs and on the same chains once through the channel. Generating values are
$\corr = 0.85$ and excess kurtosis $3.0$; a finite mixture does not contain the
generating law exactly, so what is approached is the best fit within the family.}
\label{tab:em-convergence}
\begin{tabular}{@{}lcc@{}}
\toprule
update & $\corr$ & innovation excess kurtosis \\
\midrule
$1$ & $0.4402$ & $0.078$ \\
$30$ & $0.8412$ & $1.275$ \\
$120$ & $0.8499$ & $3.019$ \\
$600$ & $0.8504$ & $2.980$ \\
\midrule
\multicolumn{3}{@{}l@{}}{\emph{where each arm finishes}} \\
clean transition pairs & $0.8509$ & $2.939$ \\
through the channel & $0.8526$ & $2.729$ \\
\bottomrule
\end{tabular}
\end{minipage}
\end{center}


\noindent The two halves say different things. The trace shows the coordinates
separating: the correlation is essentially in place by thirty updates while the
innovation shape is still climbing at a hundred and twenty, because the mixture
label is a latent variable of its own and the channel compounds the missing
information, the mechanism of \citet{dempster1977maximum} appearing in the
convergence rate, where the theory places it.

The channel has two distinct effects, and they should not be collapsed into
one. First, it changes the local rate of EM through the fraction of missing
information, which is the effect the trace above displays. Second, at finite
sample size it reduces the observed Fisher information and can therefore
increase estimation uncertainty, an effect measured directly, by coordinate,
in Section~\ref{sec:em-information}. Under identifiability and correct
specification it need not move the population optimum, and that is the limited
sense in which the clean and channel arms may approach the same endpoint while
doing so at different rates. The endpoints here are consistent with that
reading: the two arms finish within $0.002$ of each other on the correlation
and within a fifth of a unit on excess kurtosis, both close to the generating
values. This is evidence of similar fitted endpoints in this configuration, not
of an absent information or accuracy penalty.

\section{The convergence rates of the two coordinates}
\label{sec:em-two-rates}

That is one configuration. A dedicated sweep measured the cost across sixteen
seeds, recording for each fit the first update after which a coordinate stays
within a fixed relative tolerance of its value at the end of the run. At
$\nseq = \convsize$, $C = \convcomps$ and the \convdesign{} noise design, over
$\convseeds$ seeds run to $\convcap$ updates, the correlation coefficient
settles at a median of $\convrhomed$ updates and the innovation shape at
$\convshapemed$, the largest observed value being $\convshapemax$. Forming the
ratio within each fit and then taking the median over fits gives $\convratio$;
the quotient of the two medians is $\convratioofmedians$, and the paired figure
is the one reported, because both coordinates are measured on the same fit.

Two things about what that measures. ``Settled'' is trace stability, the
first update after which a coordinate stays within $10^{-3}$ of its end-of-run
value for $\alpha$, or $2\times10^{-2}$ for the excess kurtosis, and not
distance to an optimiser; a coordinate can sit still while the density it
belongs to moves along a shallow direction. And the run has to be long enough
for the statistic to mean anything: the sweep was run at two lengths, and these
medians come from the $\convcap$-update shard, the only one in which no
configuration settles near its own cap. Against them, the fixed budget of
$\convbudget$ iterations that several experiments in
Chapter~\ref{ch:em-results} used as a default
falls short of the shape coordinate's median by about a factor of six, and
leaves $\convunsettled$ of the sweep's $\convunsettledof$ configurations
unsettled in it. That count pools both shards, which is conservative:
censoring can only shorten a measured settling time, so it cannot inflate
the count.

\section{The consequence for every non-Gaussian claim}
\label{sec:em-stopping}

\begin{remark}[Fixed budgets compare convergence rates, not estimators]
\label{rmk:em-fixed-budget}
This is the single methodological lesson of the chapter, and it recurs
throughout the next one. Whatever coordinate of a fit converges slower looks
worse under a fixed iteration count, and a slower-converging configuration
under-iterated by the same budget produces what looks like a genuine deficit
rather than symmetric noise, which is precisely the signature that makes it
easy to mistake for a real effect. Compared correctly, at eight seeds and a
fixed budget of 200 updates, no channel penalty is resolved in the reported
shape summary: the paired clean $-$ noised difference at that operating point
is $-0.052\pm0.217$, consistent with zero. The operating point is not a
converged one, the chapter's own median settling time for the shape
coordinate is $\convshapemed$ updates, so this is a statement about the
paired difference at $200$ updates, not about the converged estimators. A single non-converged draw would
not have supported that comparison in either direction. Where Chapter~\ref{ch:em-results}
reports a result obtained at a fixed, possibly insufficient budget, it says so
explicitly and states which quantities are and are not affected.
\end{remark}
